\section{Methods}

\subsection{Study selection}

All \textit{Arabidopsis thaliana} studies in OSDR were retrieved through the
\texttt{biodata} v2 metadata API, returning sample-level factor values for
\NumArabidopsisStudiesScreened{} studies. Studies were admitted to a cohort only where a
numeric time axis could be parsed from a factor-value field and the assay was Illumina
RNA-seq. Cohort~A ($\gamma$-irradiation) comprises \NumCohortAStudies{} studies and
cohort~B (spaceflight development) \NumCohortBStudies{}.
\NumExcludedStudies{} studies were excluded with a recorded reason
(\texttt{scripts/cohorts.py}).

Two exclusions matter. OSD-219 re-deposits all 32 of OSD-218's BioSample identifiers
verbatim; including both would double-count every sample. This is re-detected each run
from sample-name overlap rather than hard-coded, so the exclusion cannot outlive the
condition that justifies it. The microarray irradiation studies (OSD-46, OSD-320,
OSD-329) were held out of the primary analysis because pooling them with RNA-seq would
confound platform with timepoint.

\subsection{Pseudo-time-series construction}

Raw RSEM counts and GeneLab runsheets were retrieved from the OSDR file API. Counts were
converted to CPM and $\log_2$-transformed \emph{within} each study, then expressed as a
$\log_2$ ratio against that study's own untreated controls. This within-study referencing
removes most between-study offset by construction, because a per-study additive term
cancels between numerator and denominator. A residual per-study offset was then fitted as
the median gene-wise difference from the cross-study consensus, computed only at the
\NumSeriesSharedTimepoints{} timepoints more than one study measured; a study sharing no
timepoint with any other is left unadjusted rather than shifted against a reference it
never overlaps.

Because every value is already a ratio against the untreated state, an explicit $t=0$
column of zeros is prepended and DREM is run with
\texttt{No normalization/add 0}. Omitting either would make DREM re-anchor on the 10-min
timepoint, which is already strongly induced, and the entire early response would vanish
from the model.

OSD-782 is held out rather than pooled. It is an independent laboratory \emph{and} a
different dose regime (10--100~cGy against the core's high-dose $\gamma$), so fitting it a
shared offset would remove real biology.

Batch correction was verified against six canonical DNA-damage-response genes
(BRCA1, RAD51, PARP2, SMR7, GMI1, XRI1). The pipeline fails if fewer than three remain
2-fold responsive in the wild-type arm --- the check exists to catch a correction that has
flattened the signal it was meant to preserve.

\subsection{TF--gene prior}

A binding prior was built from the \textit{Arabidopsis} DAP-seq cistrome
\citep{omalley2016cistrome}, retrieved as GEO GSE60143. Of \NumDapseqPeakFiles{} peak
files, \NumDapseqResolvedByAGI{} name the assayed TF by AGI locus and
\NumDapseqResolvedBySymbol{} by gene symbol; symbols were resolved through a map built
from the Ensembl Plants release-54 annotation, TAIR gene aliases
\citep{berardini2015tair} and ConnecTF \citep{brooks2021connectf}.
\NumDapseqUnresolved{} files remain unresolved (ChIP-seq controls, non-\textit{Arabidopsis}
orthologue assays, and post-2014 symbols) and are listed in the run report. Resolving
symbols is not cosmetic: the unresolved set is dominated by the NAC, MYB, WRKY and bZIP
families, which is where the DNA-damage regulators are.

Peaks were assigned to genes within 1~kb of a TSS using the same Ensembl release-54
TAIR10 annotation GeneLab used to produce the count matrices, so binding and expression
share one coordinate system and one locus vocabulary. Edge scores were rank-normalised
within each TF, so a score encodes binding confidence relative to that TF's other sites
rather than its sequencing depth. The prior contains \NumPriorEdges{} edges over
\NumPriorTFs{} TFs, with a median of \NumPriorMedianTargets{} targets per TF.

\paragraph{SOG1.} DAP-seq does not assay SOG1, so the DAP-seq prior cannot attribute any
bifurcation to the very regulator this cohort is about. SOG1 edges were therefore called
from the ChIP-seq of the same experiment (OSD-496 / GEO GSE112529), in which
SOG1-3xFLAG was immunoprecipitated 20~min and 1~h after irradiation. OSDR holds only raw
reads for OSD-496, so coverage was taken from GEO. A 200~bp bin was retained where IP
exceeded \emph{both} its matched input and the no-FLAG wild-type IP two-fold, which
removes chromatin-accessibility and antibody-background artefacts that either control
alone would miss. This yielded \NumSogOnePeaksTwentyMin{} and \NumSogOnePeaksOneHour{}
peaks at the two timepoints, and \NumSogOneTargets{} target genes after the same
TSS-proximity assignment used for DAP-seq.

\subsection{The auto-decoder lever}

Cell-type context comes from an auto-decoder trained on the GSE226097 single-nucleus
\textit{Arabidopsis} atlas \citep{lee2023atlas}, which provides
\NumAtlasSignatureGenes{} signature genes across \NumAtlasClusters{} clusters and a
\NumLatentDims{}-dimensional latent code per cluster.

Each timepoint's mean linear CPM profile was deconvolved against the signature matrix by
non-negative least squares, with both sides scaled to unit norm and the atlas signatures
returned to linear space beforehand. Deconvolution uses absolute expression, never the
$\log$-ratio series: a fraction is a statement about composition, and a fold-change has
already divided composition out. \NumDeconvGenesMatched{} signature genes matched, with a
median relative residual of \NumDeconvMedianResidual{}.

Writing $\bar f_c$ for the timepoint-averaged fraction of cluster $c$ and $\hat e(g,c)$
for gene $g$'s cluster profile scaled to its own maximum, the context of a gene is
$\mathrm{ctx}(g) = \sum_c \bar f_c\, \hat e(g,c)$, and each edge is re-weighted as
\[
  w(\mathrm{TF}, g) \;=\; \mathrm{binding}(\mathrm{TF}, g)\;
  \cdot\; \mathrm{ctx}(\mathrm{TF})^{\alpha}\;\cdot\;\mathrm{ctx}(g)^{\beta},
\]
with $\alpha = 1$ and $\beta = 0.5$: the TF side dominates because DREM's attribution
asks which regulator explains a split, not which target. Weights are renormalised to
$[0,1]$ within each TF and floored at $0.05$ where binding evidence exists, so context can
down-weight an edge but never silently delete it.

Genes outside the atlas signature set are, by construction, genes whose expression varies
little across the \NumAtlasClusters{} clusters; their context is therefore close to
uniform and their edges are passed through unchanged. \NumWeightedEdgesInformed{} edges
(\PctWeightedEdgesInformed{}\%) are atlas-informed and
\PctWeightedEdgesChanged{}\% move by more than $0.01$.

\subsection{Radiation signature and decoder}

Gene sets were derived from the two-genotype arms. A gene is \emph{SOG1-dependent} when it
reaches $|\log_2 \mathrm{FC}| \ge 1$ in wild type, is a SOG1 ChIP target, and its
\textit{sog1-1} amplitude is below 40\% of its wild-type amplitude --- a genetic criterion,
not a correlation with exposure. The contrast set (\emph{SOG1-independent}) holds induced
SOG1 targets that respond in both genotypes.

Each OSDR study was scored on every contrast its metadata supports: for each factor, any
level matching a control vocabulary (ground control, non-irradiated, mock, untreated, zero
dose) defines the reference and every other level of that factor is scored against it.
Factors describing what a sample \emph{is} rather than what was done to it --- organ,
tissue, age, ecotype --- are excluded. Genes are ranked by $\log_2$ fold change and a set's
mean normalised rank is compared with 2,000 size-matched random sets from the same
ranking, giving a $z$-score and an empirical $p$. Ranks, rather than fold-change
thresholds, keep the score comparable across platforms and sequencing depths.

The decision index is the conjunction $\min(z_{\text{SOG1-dep}}, -z_{\text{MYB3R-rep}})$ at
a fixed threshold of $1.96$, so both arms must be individually significant. The threshold
is statistical rather than fitted: a threshold placed between the labelled classes is
hostage to OSDR's positives including a 10~cGy exposure that produces no response.

\subsection{Tissue attribution}

Tissue association was tested against the atlas expression matrix directly, not through
the deconvolution. Each gene's cluster profile is scaled to its own maximum, a set's mean
scaled expression per cluster is compared with 1,000 size-matched random sets, and the
result is a per-cluster $z$ and empirical $p$. Deconvolution stability was measured as the
median relative change in fraction between adjacent timepoints, restricted to compartments
carrying at least 2\% mass, and the time-averaged composition was bootstrapped over
timepoints.

\subsection{Spaceflight TF activity, statistics and machine learning}

TF activity is the rank statistic above applied to each transcription factor's target set
(edges scoring $\ge 0.5$, at least 20 targets), giving \NumTFFeatures{} features per
contrast. The null is analytic rather than permuted: for a size-$n$ subset of $N$ ranks the
subset mean has variance $\frac{\sigma^2}{n}\cdot\frac{N-n}{N-1}$, the finite-population
correction, which is exact. Agreement with permutation was verified across set sizes from
20 to 3{,}000 (\texttt{--verify-null}); differences were within Monte-Carlo error.

Per-TF tests average contrasts within a mission before testing, so the unit of analysis is
the mission. Flight effects are one-sample $t$ tests over mission means; flight-versus-
radiation uses Welch's $t$; both are Benjamini--Hochberg corrected across TFs.

Classification is L2-regularised logistic regression ($C = 0.1$, balanced classes) and a
random forest, under leave-one-mission-out cross-validation, with radiation studies each
forming their own group. Significance comes from permuting the group-to-label assignment,
not the contrast labels. A platform classifier trained on identical features under
identical cross-validation serves as the confound control.

Trajectory projection correlates each contrast's TF-activity vector (Spearman) against the
same vector computed at each DREM timepoint. Because an argmax over eight profiles always
exists, the reported statistics are the shape of the mean correlation profile and the
fraction of per-contrast argmaxes falling at or beside its peak.

\subsection{Model fitting and comparison}

Models were fitted with DREM~2.0.7 \citep{dremsoftware} in batch mode
(\texttt{java -jar drem.jar -b}), with a pinned random seed, node penalty 40, at most
three paths out of a split, and a $|\log_2| \ge 1$ expression filter. Replicates were
supplied as separate full matrices rather than pre-medianed. The wrapper was validated
against DREM's own bundled example before any study data was fitted.

Each arm was fitted three times under identical settings, differing only in the prior:
\emph{flat} (all edges 1.0), \emph{binding} (DAP-seq/ChIP score), and \emph{weighted}
(cell-type-weighted). Because seed and data are held constant, any difference in TF
attribution is attributable to the prior alone.
